
\documentclass[submission,copyright,creativecommons]{eptcs}
\providecommand{\event}{ACT 2026} % Name of the event you are submitting to
\usepackage{breakurl}             % Not needed if you use pdflatex only.
\usepackage{underscore}           % Only needed if you use pdflatex.

\usepackage{amsmath,amssymb}
\usepackage{tikz-cd}

\title{Functorial Composition over Task Decomposition: \\ A Categorical Kernel for AI-Human Computation}
\author{Anonymous Author
\institute{mo:os Project\\ City, Country}
\email{author@example.com}
}
\def\titlerunning{Functorial Composition over Task Decomposition}
\def\authorrunning{A. Author}

\begin{document}
\maketitle

\begin{abstract}
Current multi-agent AI frameworks decompose tasks in a top-down manner, producing fragile orchestrations that lack formal structural guarantees. As large language models fragment tasks into sub-routines, the recombination relies on ad-hoc aggregation, brittle to hallucination and context overload. We present \texttt{mo:os}, a runtime kernel functioning as a symmetric monoidal category where all computational entities are modeled uniformly as objects (Containers). State changes are strictly defined as morphisms (ADD, LINK, MUTATE, UNLINK) over this graph topology. Applying Lawvere's functorial semantics, \texttt{mo:os} separates the syntax of the system (a JSON-schema-validated graph ontology) from its runtime semantics (a deterministic Go execution engine). By restricting LLMs to emit strictly typed morphisms rather than orchestration plans, complex behavior emerges from bottom-up functorial composition rather than unstructured decomposition. This paper introduces the category of recursive Containers, demonstrates closure under parallel and sequential wiring operations, and presents the Recursive Semantic Bridge—a functorial projection translating graph state across database, caching, and UI seamlessly without intermediate business logic. We conclude with a practical benchmark showing how this categorical architecture naturally supports multi-path (DAG) neuro-symbolic reasoning.
\end{abstract}

\section{Introduction}
The rapid advancement of generative AI has produced a crisis of orchestration. Modern frameworks such as LangChain, CrewAI, and AutoGen orchestrate complex behaviors by relying on Large Language Models (LLMs) to perform top-down task decomposition. 

\section{Background: Functorial Semantics}
Following Lawvere \cite{lawvere1963}, we establish a strict separation between syntax and semantics. Furthermore, we rely heavily on the compositional structures developed by Fong and Spivak \cite{fong2018}, treating the wiring diagrams of our container interactions as operadic compositions \cite{spivak2013}.

\section{The Container Category}
We define the Container Category $\mathcal{C}$ where objects are encapsulated nodes with typed I/O ports, and morphisms are operations upon the graph state. Formally, for any two containers $A$ and $B$, a morphism $\phi: A \to B$ represents a wiring operation or a state mutation that preserves the structural integrity defined by the global ontology.

\begin{equation}
\begin{tikzcd}
A \arrow[rr, "\phi"] \arrow[dr, "\text{add}"'] & & B \\
& C \arrow[ur, "\text{link}"'] & 
\end{tikzcd}
\end{equation}


\section{The Semantic Bottleneck: Functorial Composition vs. Task Decomposition}
The prevailing paradigm for multi-agent orchestration—exemplified by frameworks such as LangChain, CrewAI, and AutoGen—relies fundamentally on \textit{top-down task decomposition}. In these systems, a Large Language Model (LLM) is prompted to decompose a high-level objective into sub-routines, execute them iteratively, and synthesize the results. Mathematically, this relies on the semantic embedding space of the LLM to act as both the computation engine and the state router. Conducive to unbounded text injection, this lack of structural rigidity inevitably leads to context degradation, hallucination of non-existent execution paths, and brittle error recovery \cite{logicgraph2025}. We formally define this structural failure as the \textit{semantic bottleneck}.

To resolve the semantic bottleneck, \texttt{mo:os} replaces top-down decomposition with \textit{bottom-up functorial composition}. Drawing from Spivak and Rupel's formalization of temporal wiring diagrams \cite{spivak2013operad}, we restrict the LLM to function strictly as a Fuzzy Processing Unit (FPU). The FPU no longer orchestrates behavior; instead, it generates candidate morphisms (e.g., \textsc{Mutate}, \textsc{Link}, \textsc{Add}) over a strictly typed categorical state machine. The deterministic Go kernel acts as the semantic evaluator, accepting or rejecting these morphisms against the ontology prior to graph mutation.

This inversion of control yields significant mathematical guarantees. By defining agent interactions as structure-preserving functors across relational networks \cite{bonchi2017functorial}, the aggregate behavior of the system is no longer reliant on the LLM's capacity to maintain a coherent global state diagram. Instead, complex behavior emerges operadically \cite{spivak2013}. The state of the system is closed under composition: the wiring of two valid graph states invariably produces a third valid graph state, yielding a multi-agent OS that is structurally immune to orchestration drift.

\section{The Recursive Semantic Bridge}
We formalize the \texttt{mo:os} architecture as a realization of Lawvere's functorial semantics, where the system's runtime state is not a collection of ad-hoc objects, but a category-theoretical evaluation of an abstract program. In this model, the database graph serves as the \textit{Abstract Syntax Tree} (AST) of our language. We define a syntax category $\mathbf{Syn}$ where objects are Container types (Nodes) and morphisms are wiring diagrams (Edges).

The \textit{Recursive Semantic Bridge} is defined as a structure-preserving functor $F: \mathbf{Syn} \to \mathbf{Sem}$, where $\mathbf{Sem}$ is the category of deterministic Go-kernel evaluations. This functor maps the static syntax of the graph onto the dynamic rest state of the operating system.

\begin{equation}
\begin{tikzcd}
\mathbf{Syn} \arrow[r, "F"] \arrow[d, "\alpha"'] & \mathbf{Sem} \arrow[d, "F(\alpha)"] \\
\mathbf{Syn}' \arrow[r, "F"'] & \mathbf{Sem}'
\end{tikzcd}
\end{equation}

Crucially, the Go kernel acts as the semantic evaluator—accepting, validating, and committing predicted morphisms (morphisms generated by the FPU) to the persistent AST. This architecture eliminates the traditional "semantic gap" found in conventional software, where business logic must manually bridge data models to execution threads.

By treating the graph as a recursive semantic bridge, we achieve \textit{Harness Agnosticism}. Any external tool or model is merely an object in $\mathbf{Syn}$, and its interaction with the global graph is governed by the functorial mapping. This ensures that the UI (the "Surface") is a direct, consistent projection of the categorical state, removing the need for intermediary state-management synchronization. The recursive nature of the bridge allows for nested sub-workspaces, where each child workspace is a submodule in the larger categorical composition, satisfying the condition of global structural integrity via bottom-up algebraic closure.

\section{System 3: Multi-Path DAG Reasoning}
We introduce a mechanism extending LogicGraph \cite{logicgraph2025} to explore vast parallel operational states. Unlike System 1 (intuitive token prediction) and System 2 (linear chain-of-thought), System 3 reasoning treats the state space as a Directed Acyclic Graph (DAG) of potential categorical mutations. The deterministic Go kernel spawns candidate sub-graphs, where each edge represents a structure-preserving morphism.

By applying Monte Carlo Tree Search (MCTS) or similar heuristic explorations over the morphism space, \texttt{mo:os} can simulate thousands of branching future states prior to commitment. Because each branch is restricted to valid categorical syntax, the search space is pruned of logically inconsistent outcomes. This allows for high-stakes reasoning—such as complex infrastructure deployment or data migration—where the error surface of a single LLM prompt would be prohibitively large.

\section{Implementation}
\texttt{mo:os} is implemented conceptually in Go and PostgreSQL, demonstrating the math handles industrial workloads with zero data loss.

\section{Related Work}
While Catlab provides scientific libraries and Statebox \cite{statebox2019} targets blockchains, \texttt{mo:os} builds an OS constraint framework around LLM reasoning.

\section{Conclusion}
In this paper, we have demonstrated that the transition from top-down task decomposition to bottom-up functorial composition provides a necessary foundation for the next generation of multi-agent operating systems. By anchoring Large Language Models as coprocessors (FPUs) within a strictly typed categorical harness, \texttt{mo:os} resolves the semantic bottleneck, guaranteeing structural integrity and data sovereignty. Our framework proves that mathematical rigor is not an obstacle to agentic autonomy, but rather its most robust enabling technology. Future work will focus on the quantitative benchmarking of morphism-based error recovery rates against traditional LLM orchestration.

\nocite{*}
\bibliographystyle{eptcs}
\bibliography{references}

\end{document}
